% !TeX program = xelatex
% UTF-8 transcription in the order of the eight original handwritten pages.
% Original numbers, doses, abbreviations, and wording are retained.
% Source images: images/pic01.jpg through images/pic08.jpg.
% Standalone XeLaTeX document, also usable with \input in an Xepersian document.
\ifdefined\settextfont
  \def\PericardiumStandalone{0}
\else
  \def\PericardiumStandalone{1}
  \documentclass[12pt,a4paper]{article}
  \usepackage[margin=2cm]{geometry}
  \usepackage{amsmath,amssymb,enumitem}
  \usepackage{xepersian}
  \IfFontExistsTF{Amiri}{\settextfont{Amiri}}{\settextfont{FreeSerif}}
  \setlatintextfont{FreeSerif}
\fi
\providecommand{\en}[1]{\lr{#1}}
\newcommand{\peripage}[2]{\clearpage
  \section*{براونوالد ـ ۸۲ ـ بیماری‌های پریکارد ـ صفحهٔ #1}
  {\footnotesize\en{#2}\par}}
\ifnum\PericardiumStandalone=1 \begin{document}\fi

\peripage{۱}{pic01.jpg — 20261012\_183439006.jpg}
\subsection*{پریکاردیت حاد}
\begin{itemize}
\item یک سندرم التهابی \en{$\pm$} پریکاردیال افیوژن است.
\item در کشورهای در حال توسعه \en{TB} شایع‌ترین علت است.
\item ۵\% \en{CP}های غیر ایسکمیک در اورژانس، پریکاردیت هستند.
\item ایدیوپاتیک: با تست‌های روتین علتی یافت نشود؛ معمولاً ویروسی در نظر گرفته می‌شوند.
\item اگر همراه با میوکاردیت باشد: \en{Trop} تا حد متوسطی بالا می‌رود.
\item \en{LV dysfunction} نادر است.
\item \en{LVD $\ominus$}: مایوپریکاردیت.
\item \en{LVD $\oplus$}: پری‌مایوکاردیت.
\end{itemize}
\subsection*{یافته‌های بالینی}
\begin{itemize}
\item درد شدید قفسهٔ سینه در ۹۰\%؛ با دراز کشیدن بدتر می‌شود.
\item درد پلوریتیک، سوزنده، مبهم.
\item شرح حال بیماری ویروسی اخیر.
\item بیمار با پریکاردیت حاد غیر کامپلیکه: مضطرب، تب خفیف \en{$<38$}، تاکی‌کاردی سینوس.
\item آریتمی ناشایع است: \en{AF} و فلاتر در ۵\%.
\end{itemize}

\peripage{۲}{pic02.jpg — 20261012\_183445369.jpg}
\begin{itemize}
\item سمع قلب: در \en{$1/3$} موارد فریکشن راب.
\item ممکن است \en{On/off} باشد؛ به همین دلیل مرتب باید سمع نمود.
\item راب کلاسیک ۳ جزء دارد:
  \begin{itemize}
  \item سیستول بطنی.
  \item دیاستول اولیه.
  \item سیستول دهلیز.
  \end{itemize}
\item در \en{lower left sternal border} بهتر سمع می‌شود، موقعی که بیمار به جلو خم شود.
\end{itemize}
\subsection*{سایر تست‌ها}
\begin{itemize}
\item \en{EKG}: \en{STE} منتشر.
\item تغییرات نوار می‌توانند داینامیک باشند.
\item \en{PR Segment Dep}: زودرس‌ترین تغییر پریکاردیت است.
\end{itemize}
\subsection*{سیر تغییرات \en{EKG}}
\begin{enumerate}
\item \en{PRD + STE} منتشر.
\item نرمال شدن \en{ST}.
\item \en{TWI} با یا بدون \en{STD}.
\item نرمال شدن نوار.
\end{enumerate}

\peripage{۳}{pic03.jpg — 20261012\_183451228.jpg}
\begin{itemize}
\item تغییر در \en{EKG} به معنی درگیری میوکارد است، چون پریکارد از نظر الکتریکی \en{Silent} است؛ صرفاً در ۶۰\% تغییرات دارد.
\item اگر نوار \en{low voltage} بود، به معنی افیوژن شدید است.
\end{itemize}
\subsection*{آزمایش‌ها}
\begin{itemize}
\item اکثر بیماران افزایش خفیف \en{WBC} در حد \en{12--14} هزار دارند.
\item \en{CRP} در \en{$3/4$} بیماران افزایش‌یافته است.
  \begin{itemize}
  \item معمولاً ظرف ۱ هفته و حداکثر ۴ هفته \en{N/L} می‌شود.
  \item اگر \en{N/L} نشه یعنی عود آینده دارد.
  \item برای مانیتورینگ از آن استفاده می‌شود.
  \end{itemize}
\item \en{CXR} در موارد غیر کامپلیکه \en{N/L} است.
\item حتی کوچک‌ترین میزان افزایش سایز قلب، به معنی حداقل \en{300 ml} افیوژن است.
\item اکو: در ۴۰\% بیماران کاملاً \en{N/L} است.
\item در هر بیمار با شک به پریکاردیت باید انجام شود.
\item ۶۰\% بیماران افیوژن دارند:
  \begin{itemize}
  \item \en{Small ($<10$ mm)}.
  \item \en{Moderate (10--20 mm)}.
  \item \en{Large ($>20$ mm)}.
  \end{itemize}
\end{itemize}

\peripage{۴}{pic04.jpg — 20261012\_183458547.jpg}
\begin{itemize}
\item وجود افیوژن در بیماری با شرح حال منطبق بر پریکاردیت حاد، به معنی تأیید تشخیص است.
\item وقتی تشخیص قطعی یا محتمل شد: شروع درمان با \en{NSAID} + کلشی‌سین.
\item اگر تشخیص نامعلوم بود: \en{CT} یا \en{CMR}.
\end{itemize}
\subsection*{کرایتریای تشخیص بالینی: حداقل ۲ مورد}
\begin{enumerate}
\item درد قفسهٔ سینه.
\item سمع فریکشن راب.
\item تغییرات نوار: \en{STE} با یا بدون \en{PRD}.
\end{enumerate}
\begin{itemize}
\item در \en{CT}: ضخیم شدن یا هایپرانهنسمنت پریکارد.
\item در \en{CMR}: ادم پریکارد؛ \en{fat suppressed $T_2$-weighted dark blood image}.
\item روتین اولیهٔ هر بیمار با شک به پریکاردیت: \en{Cr, CRP, Trop, CXR, Echo}.
\end{itemize}

\peripage{۵}{pic05.jpg — 20261012\_183507011.jpg}
\begin{itemize}
\item شک به \en{SLE}: چک \en{ANA}.
\item اگر یافته‌ها و آزمایشات به نفع پریکاردیت حاد نبود، درمان با \en{NSAID} را شروع می‌کنیم.
\item محدودیت فعالیت فیزیکی تا زمان \en{N/L} شدن \en{CRP} و رفع علائم بالینی توصیه می‌شود.
\item ورزشکاران تا ۳ ماه نباید فعالیت ورزشی انجام دهند.
\end{itemize}
\subsection*{رژیم‌های توصیه‌شده}
\begin{enumerate}
\item \en{Tab Ibuprofen 600--800 mg PO TDS}.
\item \en{Tab ASA 750--1000 mg PO TDS}.
\end{enumerate}
\begin{itemize}
\item \en{PPI} نیز داده شود.
\item طول درمان \en{10--14} روز است، یا تا زمانی که \en{CRP}، \en{N/L} شود.
\item وقتی که علائم رفع و \en{CRP} منفی شده بود، دارو باید \en{Taper} شود.
\item کلشی‌سین به عنوان مکمل \en{NSAID} برای ۳ ماه داده می‌شود:
  \begin{itemize}
  \item \en{$>70$ kg}: \en{0.5--0.6 mg PO BD}.
  \item \en{$<70$ kg}: \en{0.5--0.6 mg PO daily}.
  \end{itemize}
\end{itemize}

\peripage{۶}{pic06.jpg — 20261012\_183514829.jpg}
\begin{itemize}
\item اگر خواستیم \en{NSAID} سریع‌تر اثر کند، می‌توانیم وریدی بدهیم.
\end{itemize}
\subsection*{پیش‌بینی‌کننده‌های پروگنوز بد}
\begin{itemize}
\item ماژور:
  \begin{itemize}
  \item \en{$T>38$}.
  \item شروع تحت حاد.
  \item افیوژن شدید.
  \item تامپوناد.
  \item عدم پاسخ به درمان طی یک هفته.
  \end{itemize}
\item مینور:
  \begin{itemize}
  \item مایوپریکاردیت.
  \item ایمونوساپرشن.
  \item تروما.
  \item آنتی‌کوآگولان خوراکی.
  \end{itemize}
\end{itemize}
\subsection*{تأیید تشخیص پریکاردیت حاد}
\begin{itemize}
\item آیا شک قوی به اتیولوژی خاصی وجود دارد؟ / آیا پیش‌بینی‌کننده‌هایی برای پروگنوز بد وجود دارد؟
  \begin{itemize}
  \item بله: بیمار \en{High Risk}؛ بستری.
  \item خیر: شروع درمان امپریکال.
  \end{itemize}
\item آیا به درمان جواب داد؟
  \begin{itemize}
  \item بله: درمان سرپایی.
  \item خیر: ریسک متوسط؛ بستری.
  \end{itemize}
\end{itemize}

\peripage{۷}{pic07.jpg — 20261012\_183522053.jpg}
\subsection*{نحوهٔ \en{Taper} کردن}
\begin{itemize}
\item آسپرین: به صورت هفتگی کاهش دوز / قطع پس از \en{2--3} هفته.
\item ایبوپروفن: به صورت هفتگی کاهش دوز / قطع پس از \en{2--3} هفته.
% Ditto marks for ibuprofen in the source repeat the aspirin instructions.
\item کلشی‌سین: طی \en{2--3} هفته.
\end{itemize}
\subsection*{کورتیکواستروئید}
\begin{itemize}
\item نباید استفاده شوند مگر اندیکاسیون خاص.
\item خطر عود را بالا می‌برند و روند درمان را کند می‌کنند.
\end{itemize}
\subsection*{اندیکاسیون کورتیکواستروئید}
\begin{enumerate}
\item کنتراندیکاسیون \en{NSAID}، عدم پاسخ به \en{NSAID}.
\item وجود بیماری زمینه‌ای که درمان آن کورتیکواستروئید است.
\item بیماری همراه (\en{AKI / CKD}).
\item بارداری.
\item داروهایی که با \en{NSAID} یا کلشی‌سین تداخل دارند: آنتی‌کوآگولان‌های خوراکی.
\end{enumerate}
\begin{itemize}
\item دوز کم آن توصیه می‌شود: \en{Prednisone 0.2--0.5 mg/kg daily}.
\item کورتیکواستروئید باید ظرف \en{9--12 w} تیپر شود.
\item کلشی‌سین هم باید بدهیم.
\end{itemize}

\peripage{۸}{pic08.jpg — 20261012\_183527489.jpg}
\subsection*{پریکاردیت راجعه (ریکارنت)}
\begin{itemize}
\item در \en{15--30\%} بیماران با پریکاردیت حاد ایدیوپاتیک.
\item بین دو اپیزود باید حداقل \en{4--6} هفته فاصله باشد.
\item درمان تفاوتی با فرم اولیه ندارد.
\item در عود سوم \en{Triple} تراپی می‌کنیم:
  \begin{itemize}
  \item \en{NSAID} + کورتون + کلشی‌سین.
  \item \en{Prednisone 0.2--0.5 mg/kg/day for 2--4 w}؛ تیپر طی \en{2--4} هفته.
  \end{itemize}
\item کلشی‌سین برای ۶ ماه می‌دهیم؛ تا ۱۲ ماه.
\end{itemize}
% Each original page repeats the subject header and has an unfilled date field.
% Printed form labels: موضوع، تاریخ. Handwritten page numbers: ۱ through ۸.
% The small printed notebook logo contains no further medical text.
\ifnum\PericardiumStandalone=1
  \let\PericardiumEnd\enddocument
\else
  \let\PericardiumEnd\relax
\fi
\PericardiumEnd
